\chapter{Geometry of Interaction and Partial Linear Logic for Mealy Processes}
\label{ch:goi-partial-linear-logic-mealy}

\section{Overview}
\label{sec:goi-overview}

The previous chapter showed that the symmetric monoidal category
\[
\mathsf{Mly}(\mathcal E)
\]
of internal categories and Mealy morphisms carries a canonical partial trace.
This trace is inherited from the total compact-closed trace on spans of
internal categories.  It is defined precisely for those Mealy morphisms whose
closed-loop span remains input-discrete.  Equivalently, in component form, the
trace is defined precisely when the internal feedback equation
\[
h=\omega_U((a,h),x)
\]
has a unique solution over every visible input and closed state.

This chapter applies that partial trace to Geometry of Interaction.  The
guiding principle is that a bidirectional process from an interface
\(\mathbb A\) to an interface \(\mathbb B\) is a Mealy morphism
\[
\mathbb A^+\times\mathbb B^-
\rightsquigarrow
\mathbb B^+\times\mathbb A^-.
\]
The positive component flows forward, while the negative component flows
backward.  Composition is not ordinary sequential composition of the displayed
Mealy morphisms.  Instead, composition is Geometry-of-Interaction execution:
one first forms an open pipeline and then closes the shared counterflow
interface by the partial trace.

Thus, for bidirectional Mealy processes
\[
F:\mathbb A\to\mathbb B,
\qquad
G:\mathbb B\to\mathbb C,
\]
their composite is defined by
\[
G\circ_{\partial}F
=
\operatorname{Tr}^{\mathbb B^-}(H_{G,F}),
\]
where \(H_{G,F}\) is the open pipeline from \(\mathbb A\) to \(\mathbb C\)
with the \(\mathbb B^-\)-interface still exposed.  The composite is defined
exactly when this open pipeline is trace-class.

In component notation, suppose
\[
F=(P_F,\delta_F,\omega_F):
\mathbb A^+\times\mathbb B^-
\rightsquigarrow
\mathbb B^+\times\mathbb A^-,
\]
and
\[
G=(P_G,\delta_G,\omega_G):
\mathbb B^+\times\mathbb C^-
\rightsquigarrow
\mathbb C^+\times\mathbb B^-.
\]
For a visible input \((a^+,c^-)\) and a compatible closed pipeline state
\((x,y)\), the hidden counterflow witness \(h\in B^-_1\) must satisfy
\[
h
=
\omega_G^-
\left(
\left(
\omega_F^+((a^+,h),x),
c^-
\right),
y
\right).
\]
When this equation has a unique solution, denoted
\[
h_{G,F}(a^+,c^-,x,y),
\]
the composite process is obtained by substituting that solution into the
transitions and outputs of \(F\) and \(G\).

The resulting structure is a symmetric monoidal paracategory
\[
\operatorname{Int}_{\partial}(\mathsf{Mly}(\mathcal E)).
\]
Its objects are bidirectional Mealy interfaces, its morphisms are
bidirectional Mealy processes, and its partial composition is GoI execution by
trace-class feedback.

This paracategory gives the process semantics of multiplicative linear logic
with partial cut.  Formulas are interpreted as bidirectional Mealy interfaces,
proofs as bidirectional Mealy processes, tensor as parallel composition,
linear negation as reversal of polarity, and cut as partial GoI execution.
A cut is defined precisely when the associated feedback equation is
trace-class.

The central slogans are:
\[
\boxed{
\text{GoI composition}
=
\text{open pipeline}
+
\text{trace-class feedback}.
}
\]
and
\[
\boxed{
\text{cut admissibility}
=
\text{unique solvability of the feedback equation}.
}
\]

\section{Standing assumptions and conventions}
\label{sec:goi-standing-assumptions}

Throughout this chapter, assume that \(\mathcal E\) has finite limits.  Then
the category
\[
\mathbf{Cat}_{\mathrm{int}}(\mathcal E)
\]
of internal categories and internal functors has finite limits, computed
componentwise.

By the previous chapter, the symmetric monoidal category
\[
\mathsf{Mly}(\mathcal E)
\]
of internal categories and Mealy morphisms carries a canonical partial trace.
The trace of a Mealy morphism is defined exactly when the corresponding
ambient traced span of internal categories remains input-discrete.

All equalities of products, spans, Mealy morphisms, traces, and composites are
read up to the canonical associativity, unit, symmetry, pullback, and span-apex
isomorphisms.  Equivalently, one may work in the quotient where such canonical
isomorphisms are identified.

We use the term \emph{paracategory} for a category-like structure with partial
composition.  In the present case, a composite is defined exactly when the
corresponding GoI feedback problem is trace-class.  Associativity and unit laws
are read on their common trace-class domains and are inherited from the trace
identities of the previous chapter.

\section{Bidirectional Mealy interfaces}
\label{sec:bidirectional-mealy-interfaces}

\begin{definition}[Bidirectional Mealy interface]
\label{def:bidirectional-mealy-interface}
A \textbf{bidirectional Mealy interface} is a pair
\[
\mathbb A=(\mathbb A^+,\mathbb A^-)
\]
of internal categories in \(\mathcal E\).
\end{definition}

The component \(\mathbb A^+\) is the positive or forward interface.  The
component \(\mathbb A^-\) is the negative or counterflow interface.

\begin{definition}[Dual interface]
\label{def:bidirectional-dual-interface}
The \textbf{dual} of a bidirectional interface
\[
\mathbb A=(\mathbb A^+,\mathbb A^-)
\]
is
\[
\mathbb A^\bot:=(\mathbb A^-,\mathbb A^+).
\]
\end{definition}

Thus there is a strict equality
\[
(\mathbb A^\bot)^\bot=\mathbb A.
\]

\begin{definition}[Tensor of interfaces]
\label{def:tensor-bidirectional-interfaces}
The tensor product of bidirectional interfaces is defined componentwise:
\[
\mathbb A\otimes\mathbb B
:=
(\mathbb A^+\times\mathbb B^+,\,
 \mathbb A^-\times\mathbb B^-).
\]
The tensor unit is
\[
\mathbf I:=(\mathbf 1,\mathbf 1),
\]
where \(\mathbf 1\) is the terminal internal category.
\end{definition}

The associativity, unit, and symmetry constraints for \(\otimes\) are inherited
from products of internal categories.

\begin{definition}[Par of interfaces]
\label{def:par-bidirectional-interfaces}
The par interface is defined by De Morgan duality:
\[
\mathbb A\parr\mathbb B
:=
(\mathbb A^\bot\otimes\mathbb B^\bot)^\bot.
\]
The par unit is
\[
\bot:=\mathbf I^\bot.
\]
\end{definition}

Since
\[
\mathbb A^\bot\otimes\mathbb B^\bot
=
(\mathbb A^-\times\mathbb B^-,\,
 \mathbb A^+\times\mathbb B^+),
\]
one has a canonical isomorphism
\[
\mathbb A\parr\mathbb B
\cong
(\mathbb A^+\times\mathbb B^+,\,
 \mathbb A^-\times\mathbb B^-).
\]
Thus tensor and par have canonically isomorphic underlying bidirectional
interfaces.  Their logical distinction is supplied by their polarity and by
the proof rules in which they occur.

\begin{proposition}[Interface-level linear structure]
\label{prop:interface-level-linear-structure}
Bidirectional Mealy interfaces carry a symmetric monoidal structure
\[
(\otimes,\mathbf I)
\]
with involutive duality
\[
(-)^\bot.
\]
Moreover,
\[
(\mathbb A\otimes\mathbb B)^\bot
\cong
\mathbb A^\bot\parr\mathbb B^\bot,
\qquad
(\mathbb A\parr\mathbb B)^\bot
\cong
\mathbb A^\bot\otimes\mathbb B^\bot,
\]
and
\[
\mathbf I^\bot=\bot,
\qquad
\bot^\bot=\mathbf I.
\]
\end{proposition}

\begin{proof}
All claims follow immediately from the definitions and from the associativity,
unit, and symmetry isomorphisms for products of internal categories.
\end{proof}

\section{Bidirectional Mealy processes}
\label{sec:bidirectional-mealy-processes}

\begin{definition}[Bidirectional Mealy process]
\label{def:bidirectional-mealy-process}
Let
\[
\mathbb A=(\mathbb A^+,\mathbb A^-),
\qquad
\mathbb B=(\mathbb B^+,\mathbb B^-)
\]
be bidirectional Mealy interfaces.  A \textbf{bidirectional Mealy process}
\[
F:\mathbb A\to\mathbb B
\]
is a Mealy morphism
\[
F:
\mathbb A^+\times\mathbb B^-
\rightsquigarrow
\mathbb B^+\times\mathbb A^-.
\]
\end{definition}

Thus a process from \(\mathbb A\) to \(\mathbb B\) consumes positive input
from \(\mathbb A^+\) and negative counterflow input from \(\mathbb B^-\), and
produces positive output in \(\mathbb B^+\) and negative counterflow output in
\(\mathbb A^-\).

Let
\[
F=(P_F,\delta_F,\omega_F):
\mathbb A^+\times\mathbb B^-
\rightsquigarrow
\mathbb B^+\times\mathbb A^-.
\]
The carrier span has the form
\[
\begin{tikzcd}[column sep=large]
A_0^+\times B_0^-
&
P_F
  \arrow[l, "p_F"']
  \arrow[r, "q_F"]
&
B_0^+\times A_0^- .
\end{tikzcd}
\]
Write
\[
p_F=(p_F^+,p_F^-),
\qquad
q_F=(q_F^+,q_F^-).
\]
The output map decomposes as
\[
\omega_F=(\omega_F^+,\omega_F^-).
\]

Given an admissible input arrow
\[
(a^+,b^-)
\]
at a state \(x\in P_F\), the process produces a new state
\[
\delta_F((a^+,b^-),x)
\]
and output arrows
\[
\omega_F^+((a^+,b^-),x)\in B_1^+,
\qquad
\omega_F^-((a^+,b^-),x)\in A_1^-.
\]

The Mealy typing equations say
\[
p_F\delta_F((a^+,b^-),x)
=
(s_{A^+}a^+,s_{B^-}b^-),
\]
and
\[
s_{B^+\times A^-}\omega_F((a^+,b^-),x)
=
q_F\delta_F((a^+,b^-),x),
\]
\[
t_{B^+\times A^-}\omega_F((a^+,b^-),x)
=
q_F(x).
\]
Equivalently, componentwise,
\[
s_{B^+}\omega_F^+((a^+,b^-),x)
=
q_F^+\delta_F((a^+,b^-),x),
\]
\[
t_{B^+}\omega_F^+((a^+,b^-),x)
=
q_F^+(x),
\]
and
\[
s_{A^-}\omega_F^-((a^+,b^-),x)
=
q_F^-\delta_F((a^+,b^-),x),
\]
\[
t_{A^-}\omega_F^-((a^+,b^-),x)
=
q_F^-(x).
\]

\section{Identity and tensor of processes}
\label{sec:identity-tensor-processes}

\begin{definition}[Identity process]
\label{def:identity-process}
For a bidirectional interface
\[
\mathbb A=(\mathbb A^+,\mathbb A^-),
\]
the identity process
\[
1_{\mathbb A}:\mathbb A\to\mathbb A
\]
is the identity Mealy morphism
\[
1_{\mathbb A^+\times\mathbb A^-}:
\mathbb A^+\times\mathbb A^-
\rightsquigarrow
\mathbb A^+\times\mathbb A^-.
\]
\end{definition}

Operationally, \(1_{\mathbb A}\) is the wire process: it passes positive
behaviour forward and negative behaviour backward.

\begin{definition}[Tensor of processes]
\label{def:tensor-processes}
Let
\[
F:\mathbb A\to\mathbb B
\]
and
\[
G:\mathbb C\to\mathbb D
\]
be bidirectional Mealy processes.  Their tensor product
\[
F\otimes G:
\mathbb A\otimes\mathbb C
\to
\mathbb B\otimes\mathbb D
\]
is the process induced by the product Mealy morphism
\[
F\times G
\]
after the canonical associativity and symmetry identifications
\[
(\mathbb A^+\times\mathbb C^+)
\times
(\mathbb B^-\times\mathbb D^-)
\cong
(\mathbb A^+\times\mathbb B^-)
\times
(\mathbb C^+\times\mathbb D^-),
\]
and
\[
(\mathbb B^+\times\mathbb A^-)
\times
(\mathbb D^+\times\mathbb C^-)
\cong
(\mathbb B^+\times\mathbb D^+)
\times
(\mathbb A^-\times\mathbb C^-).
\]
\end{definition}

\begin{proposition}[Tensor of processes]
\label{prop:tensor-processes-well-typed}
The tensor of processes is well typed.  Together with the wire processes, it
makes bidirectional Mealy interfaces and processes into a symmetric monoidal
partial graph before composition is specified.
\end{proposition}

\begin{proof}
The tensor product of Mealy morphisms is defined componentwise in
\(\mathsf{Mly}(\mathcal E)\).  Products of internal categories are computed
componentwise, and the Mealy unit and multiplication laws hold componentwise.
The required source and target objects are related by canonical product
associativity and symmetry isomorphisms.  Therefore \(F\otimes G\) has the
stated type.  The associativity, unit, and symmetry constraints are inherited
from the symmetric monoidal structure on \(\mathsf{Mly}(\mathcal E)\).
\end{proof}

\section{Open pipelines}
\label{sec:open-pipelines}

Let
\[
F:\mathbb A\to\mathbb B
\]
and
\[
G:\mathbb B\to\mathbb C
\]
be bidirectional Mealy processes.  Thus
\[
F:
\mathbb A^+\times\mathbb B^-
\rightsquigarrow
\mathbb B^+\times\mathbb A^-,
\]
and
\[
G:
\mathbb B^+\times\mathbb C^-
\rightsquigarrow
\mathbb C^+\times\mathbb B^-.
\]

The open pipeline connects the positive \(\mathbb B^+\)-wire from \(F\) to
\(G\), while leaving the negative \(\mathbb B^-\)-wire exposed.  This exposed
negative wire is later closed by the partial trace.

\begin{definition}[Open pipeline]
\label{def:open-pipeline}
The \textbf{open pipeline} of
\[
F:\mathbb A\to\mathbb B
\]
and
\[
G:\mathbb B\to\mathbb C
\]
is the Mealy morphism
\[
H_{G,F}:
\mathbb A^+\times\mathbb C^-\times\mathbb B^-
\rightsquigarrow
\mathbb C^+\times\mathbb A^-\times\mathbb B^-,
\]
obtained by first running \(F\) on the \(\mathbb A^+\)-input and exposed
\(\mathbb B^-\)-input, then running \(G\) on the \(\mathbb B^+\)-output of
\(F\) and the \(\mathbb C^-\)-input.
\end{definition}

Categorically, \(H_{G,F}\) is the ordinary Mealy composite formed by connecting
the visible \(\mathbb B^+\)-component, after the canonical product
reassociations.  The \(\mathbb B^-\)-component is not composed away; it is
kept as an exposed hidden interface for the later trace.

We now spell this out in component notation.  Let
\[
F=(P_F,\delta_F,\omega_F),
\qquad
G=(P_G,\delta_G,\omega_G).
\]
Write
\[
p_F=(p_F^+,p_F^-):P_F\to A_0^+\times B_0^-,
\]
\[
q_F=(q_F^+,q_F^-):P_F\to B_0^+\times A_0^-,
\]
and
\[
p_G=(p_G^+,p_G^-):P_G\to B_0^+\times C_0^-,
\]
\[
q_G=(q_G^+,q_G^-):P_G\to C_0^+\times B_0^-.
\]

The carrier of the open pipeline is
\[
P_{G,F}
:=
P_F\times_{q_F^+,B_0^+,p_G^+}P_G.
\]
A state is a pair
\[
(x,y)
\]
with
\[
q_F^+(x)=p_G^+(y).
\]

The source map of \(H_{G,F}\) is
\[
p_{G,F}:P_{G,F}\to A_0^+\times C_0^-\times B_0^-,
\]
given by
\[
p_{G,F}(x,y)
=
(p_F^+(x),p_G^-(y),p_F^-(x)).
\]
The target map is
\[
q_{G,F}:P_{G,F}\to C_0^+\times A_0^-\times B_0^-,
\]
given by
\[
q_{G,F}(x,y)
=
(q_G^+(y),q_F^-(x),q_G^-(y)).
\]

Given an admissible input arrow
\[
(a^+,c^-,h)
\]
at a state \((x,y)\), first compute
\[
x'
=
\delta_F((a^+,h),x),
\]
and
\[
(b^+,a^-)
=
\omega_F((a^+,h),x).
\]
Then compute
\[
y'
=
\delta_G((b^+,c^-),y),
\]
and
\[
(c^+,\bar h)
=
\omega_G((b^+,c^-),y).
\]
The open-pipeline transition and output are
\[
\delta_{G,F}((a^+,c^-,h),(x,y))
=
(x',y'),
\]
and
\[
\omega_{G,F}((a^+,c^-,h),(x,y))
=
(c^+,a^-,\bar h).
\]

\begin{proposition}[Well-definedness of open pipelines]
\label{prop:open-pipeline-well-defined}
The open pipeline \(H_{G,F}\) is a well-defined Mealy morphism
\[
\mathbb A^+\times\mathbb C^-\times\mathbb B^-
\rightsquigarrow
\mathbb C^+\times\mathbb A^-\times\mathbb B^-.
\]
\end{proposition}

\begin{proof}
The output arrow \(b^+\) of \(F\) has target
\[
t_{B^+}(b^+)=q_F^+(x)=p_G^+(y),
\]
so it is admissible as the \(\mathbb B^+\)-input to \(G\) at state \(y\).
Moreover,
\[
s_{B^+}(b^+)=q_F^+(x'),
\]
while the Mealy typing equation for \(G\) gives
\[
p_G^+(y')=s_{B^+}(b^+).
\]
Thus
\[
q_F^+(x')=p_G^+(y'),
\]
so the pair \((x',y')\) is again a state of \(P_{G,F}\).

The source and target equations for the output
\[
(c^+,a^-,\bar h)
\]
follow componentwise from the Mealy typing equations for \(F\) and \(G\).  The
unit and multiplication laws for \(H_{G,F}\) follow from the unit and
multiplication laws for \(F\) and \(G\), together with the compatibility of the
output \(b^+\) of \(F\) with the input of \(G\).  Thus \(H_{G,F}\) is a Mealy
morphism of the displayed type.
\end{proof}

\section{Partial GoI composition}
\label{sec:partial-goi-composition}

The open pipeline still exposes the negative interface \(\mathbb B^-\).  The
Geometry-of-Interaction composite is obtained by closing this exposed
interface using the partial trace of the previous chapter.

\begin{definition}[Partial GoI composition]
\label{def:partial-goi-composition}
Let
\[
F:\mathbb A\to\mathbb B,
\qquad
G:\mathbb B\to\mathbb C
\]
be bidirectional Mealy processes.  Their partial GoI composite is
\[
G\circ_{\partial}F
:=
\operatorname{Tr}^{\mathbb B^-}(H_{G,F}),
\]
whenever the open pipeline
\[
H_{G,F}:
\mathbb A^+\times\mathbb C^-\times\mathbb B^-
\rightsquigarrow
\mathbb C^+\times\mathbb A^-\times\mathbb B^-
\]
is trace-class over \(\mathbb B^-\).
\end{definition}

When defined, the result has type
\[
G\circ_{\partial}F:
\mathbb A^+\times\mathbb C^-
\rightsquigarrow
\mathbb C^+\times\mathbb A^-,
\]
and hence is a bidirectional Mealy process
\[
G\circ_{\partial}F:\mathbb A\to\mathbb C.
\]

Thus
\[
\boxed{
G\circ_{\partial}F
=
\operatorname{Tr}^{\mathbb B^-}(H_{G,F}).
}
\]

Composition is partial because the exposed counterflow feedback equation may
fail to have a unique internal solution.

\section{Execution witnesses and the feedback equation}
\label{sec:execution-witnesses-feedback-equation}

We now unpack the trace-class condition for the open pipeline.

Let
\[
F:\mathbb A\to\mathbb B
\]
and
\[
G:\mathbb B\to\mathbb C
\]
be as above.  The open pipeline has carrier
\[
P_{G,F}=P_F\times_{B_0^+}P_G.
\]
The hidden input object for the trace is the \(\mathbb B^-\)-component of the
source:
\[
p_F^-:P_F\to B_0^-.
\]
The hidden output object for the trace is the \(\mathbb B^-\)-component of the
target:
\[
q_G^-:P_G\to B_0^-.
\]

\begin{definition}[Closed pipeline state object]
\label{def:closed-pipeline-state-object}
The \textbf{closed pipeline state object} is
\[
P_{G,F}^{\circ}
:=
\operatorname{Eq}
\left(
p_F^-\pi_F,\,
q_G^-\pi_G
\right)
\subseteq
P_F\times_{B_0^+}P_G.
\]
A generalized element is a compatible pair \((x,y)\) such that
\[
q_F^+(x)=p_G^+(y)
\]
and
\[
p_F^-(x)=q_G^-(y).
\]
\end{definition}

Let
\[
D_{G,F}
:=
(A_1^+\times C_1^-)
\times_{t_{A^+}\times t_{C^-},\,
        A_0^+\times C_0^-,\,
        p_F^+\times p_G^-}
P_{G,F}^{\circ}.
\]
A generalized element of \(D_{G,F}\) is a visible input and closed state
\[
(a^+,c^-,x,y)
\]
with
\[
t_{A^+}(a^+)=p_F^+(x),
\qquad
t_{C^-}(c^-)=p_G^-(y),
\]
and \((x,y)\in P_{G,F}^{\circ}\).

Define
\[
H_{G,F}^{\mathrm{wit}}
:=
D_{G,F}
\times_{p_F^-\pi_F,\,B_0^-,\,t_{B^-}}
B_1^-.
\]
A generalized element is a tuple
\[
(a^+,c^-,x,y,h)
\]
such that
\[
(a^+,c^-,x,y)\in D_{G,F}
\]
and
\[
t_{B^-}(h)=p_F^-(x)=q_G^-(y).
\]

There are two maps
\[
H_{G,F}^{\mathrm{wit}}\to B_1^-.
\]
The first is projection:
\[
\pi_{B^-}(a^+,c^-,x,y,h)=h.
\]
The second is the hidden output produced by the open pipeline.  Define
\[
b^+
:=
\omega_F^+((a^+,h),x).
\]
Then set
\[
\chi_{G,F}(a^+,c^-,x,y,h)
:=
\omega_G^-((b^+,c^-),y).
\]
Equivalently,
\[
\chi_{G,F}(a^+,c^-,x,y,h)
=
\omega_G^-
\left(
\left(
\omega_F^+((a^+,h),x),
c^-
\right),
y
\right).
\]

\begin{definition}[Execution witness object]
\label{def:execution-witness-object}
The \textbf{execution witness object} for the pair \((G,F)\) is the equalizer
\[
W_{G,F}
:=
\operatorname{Eq}
\left(
\pi_{B^-},\chi_{G,F}
\right)
\hookrightarrow
H_{G,F}^{\mathrm{wit}}.
\]
Thus a generalized element of \(W_{G,F}\) is a tuple
\[
(a^+,c^-,x,y,h)
\]
satisfying the feedback equation
\[
h
=
\omega_G^-
\left(
\left(
\omega_F^+((a^+,h),x),
c^-
\right),
y
\right).
\]
There is a canonical projection
\[
\pi_{G,F}:W_{G,F}\to D_{G,F}.
\]
\end{definition}

\begin{theorem}[Execution criterion]
\label{thm:goi-execution-criterion}
For bidirectional Mealy processes
\[
F:\mathbb A\to\mathbb B
\]
and
\[
G:\mathbb B\to\mathbb C,
\]
the composite
\[
G\circ_{\partial}F
\]
is defined if and only if the projection
\[
\pi_{G,F}:W_{G,F}\to D_{G,F}
\]
is an isomorphism.
\end{theorem}

\begin{proof}
By definition,
\[
G\circ_{\partial}F
=
\operatorname{Tr}^{\mathbb B^-}(H_{G,F})
\]
is defined exactly when the open pipeline \(H_{G,F}\) is trace-class over
\(\mathbb B^-\).  By the internal trace-class criterion from the previous
chapter, trace-class is equivalent to the assertion that the feedback witness
projection for \(H_{G,F}\) over the visible input object is an isomorphism.

The object \(D_{G,F}\) is precisely the visible input and closed state object
for the trace of \(H_{G,F}\), while \(W_{G,F}\) is precisely the object of
hidden feedback witnesses satisfying the equation
\[
h=\chi_{G,F}(a^+,c^-,x,y,h).
\]
Thus the trace-class comparison map for \(H_{G,F}\) is exactly
\[
\pi_{G,F}:W_{G,F}\to D_{G,F}.
\]
Therefore \(G\circ_{\partial}F\) is defined if and only if
\(\pi_{G,F}\) is an isomorphism.
\end{proof}

When \(\pi_{G,F}\) is an isomorphism, denote the unique solution of the
feedback equation by
\[
h_{G,F}(a^+,c^-,x,y).
\]

\section{Component execution formula}
\label{sec:component-execution-formula}

\begin{theorem}[Component execution formula]
\label{thm:component-execution-formula}
Suppose
\[
F:\mathbb A\to\mathbb B
\]
and
\[
G:\mathbb B\to\mathbb C
\]
are bidirectional Mealy processes and that
\[
G\circ_{\partial}F
\]
is defined.  Let
\[
h_{G,F}(a^+,c^-,x,y)
\]
be the unique solution of
\[
h
=
\omega_G^-
\left(
\left(
\omega_F^+((a^+,h),x),
c^-
\right),
y
\right).
\]
Then the composite
\[
G\circ_{\partial}F:\mathbb A\to\mathbb C
\]
has carrier
\[
P_{G,F}^{\circ}
=
\operatorname{Eq}(p_F^-\pi_F,q_G^-\pi_G)
\subseteq
P_F\times_{B_0^+}P_G.
\]
For a visible input
\[
(a^+,c^-)
\]
at a closed state \((x,y)\), set
\[
h:=h_{G,F}(a^+,c^-,x,y),
\]
\[
x':=\delta_F((a^+,h),x),
\]
\[
(b^+,a^-):=\omega_F((a^+,h),x),
\]
\[
y':=\delta_G((b^+,c^-),y),
\]
and
\[
(c^+,\bar h):=\omega_G((b^+,c^-),y).
\]
Then
\[
\delta_{G\circ_{\partial}F}((a^+,c^-),(x,y))
=
(x',y'),
\]
and
\[
\omega_{G\circ_{\partial}F}((a^+,c^-),(x,y))
=
(c^+,a^-).
\]
\end{theorem}

\begin{proof}
The composite
\[
G\circ_{\partial}F
\]
is the partial trace of the open pipeline
\[
H_{G,F}
\]
over \(\mathbb B^-\).  By the component formula for partial trace in the
previous chapter, the carrier of the trace is the closed state object for the
hidden interface, namely
\[
P_{G,F}^{\circ}.
\]
The transition and output of the traced Mealy morphism are obtained by
substituting the unique hidden feedback solution into the transition and
output of \(H_{G,F}\).

The transition and output of \(H_{G,F}\) are, by definition,
\[
(x,y)
\longmapsto
\left(
\delta_F((a^+,h),x),
\delta_G((b^+,c^-),y)
\right),
\]
with
\[
b^+=\omega_F^+((a^+,h),x),
\]
and output
\[
(c^+,a^-,\bar h),
\]
where
\[
(c^+,\bar h)=\omega_G((b^+,c^-),y),
\qquad
a^-=\omega_F^-((a^+,h),x).
\]
The hidden output \(\bar h\) is discarded by the trace after being identified
with the hidden input \(h\).  Hence the visible output is
\[
(c^+,a^-).
\]
This gives the displayed formula.
\end{proof}

\section{The partial Int paracategory}
\label{sec:partial-int-paracategory}

\begin{definition}[The partial Int paracategory of Mealy processes]
\label{def:partial-int-mly}
Define
\[
\operatorname{Int}_{\partial}(\mathsf{Mly}(\mathcal E))
\]
as follows:
\begin{enumerate}
\item objects are bidirectional Mealy interfaces
\[
\mathbb A=(\mathbb A^+,\mathbb A^-);
\]
\item morphisms
\[
F:\mathbb A\to\mathbb B
\]
are bidirectional Mealy processes, equivalently Mealy morphisms
\[
F:\mathbb A^+\times\mathbb B^-
\rightsquigarrow
\mathbb B^+\times\mathbb A^-;
\]
\item identities are the wire processes
\[
1_{\mathbb A}:=1_{\mathbb A^+\times\mathbb A^-};
\]
\item tensor is defined componentwise on interfaces and by product on
processes;
\item composition is the partial GoI composition
\[
G\circ_{\partial}F
=
\operatorname{Tr}^{\mathbb B^-}(H_{G,F})
\]
when the open pipeline \(H_{G,F}\) is trace-class.
\end{enumerate}
\end{definition}

\begin{theorem}[Partial Int paracategory]
\label{thm:partial-int-paracategory}
The structure
\[
\operatorname{Int}_{\partial}(\mathsf{Mly}(\mathcal E))
\]
is a symmetric monoidal paracategory with involutive duality.  Its partial
composition is Geometry-of-Interaction execution by trace-class feedback.
\end{theorem}

\begin{proof}
The previous chapter proves that
\[
\mathsf{Mly}(\mathcal E)
\]
is a partially traced symmetric monoidal category.  Applying the partial Int
construction gives the displayed objects, morphisms, identities, tensor, and
partial composition.

The identity laws follow from the yanking law for the inherited partial trace.
Composing a process with a wire process forms an open pipeline whose feedback
equation is the wire equation, and the trace reduces to the original process.

Associativity follows from staged and simultaneous feedback.  Given three
processes
\[
F:\mathbb A\to\mathbb B,
\qquad
G:\mathbb B\to\mathbb C,
\qquad
H:\mathbb C\to\mathbb D,
\]
both bracketings
\[
H\circ_{\partial}(G\circ_{\partial}F)
\]
and
\[
(H\circ_{\partial}G)\circ_{\partial}F
\]
are obtained from the same three-stage open pipeline by tracing the same
internal counterflow interfaces in different orders.  The staged-feedback
theorem from the previous chapter identifies these iterated traces with the
simultaneous trace whenever the relevant trace-class conditions hold.

Tensor compatibility follows from the superposing law for the partial trace.
The open pipeline of
\[
F_1\otimes F_2
\]
and
\[
G_1\otimes G_2
\]
is canonically the tensor of the open pipelines of \((F_1,G_1)\) and
\((F_2,G_2)\), and tracing the tensor is identified with tensoring the traces
on the common trace-class domain.

The involutive duality is given by
\[
(\mathbb A^+,\mathbb A^-)^\bot=(\mathbb A^-,\mathbb A^+).
\]
Thus
\[
\operatorname{Int}_{\partial}(\mathsf{Mly}(\mathcal E))
\]
is a symmetric monoidal paracategory with involutive duality.
\end{proof}

\section{Linear structure of the partial Int paracategory}
\label{sec:linear-structure-partial-int}

The partial Int paracategory supports the multiplicative linear connectives.

The tensor is
\[
\mathbb A\otimes\mathbb B
=
(\mathbb A^+\times\mathbb B^+,\,
 \mathbb A^-\times\mathbb B^-),
\]
with unit
\[
\mathbf I=(\mathbf 1,\mathbf 1).
\]
Linear negation is the involution
\[
\mathbb A^\bot=(\mathbb A^-,\mathbb A^+).
\]
Par is defined by De Morgan duality:
\[
\mathbb A\parr\mathbb B
=
(\mathbb A^\bot\otimes\mathbb B^\bot)^\bot,
\]
and its unit is
\[
\bot=\mathbf I^\bot.
\]

For processes, tensor is parallel execution and cut is partial GoI
composition.  Thus the partial Int paracategory gives a process semantics for
the multiplicative fragment of linear logic in which cut is a partial
operation.

\begin{proposition}[Bidirectional currying]
\label{prop:bidirectional-currying}
For bidirectional interfaces
\[
\mathbb A,\mathbb B,\mathbb C,
\]
there is a canonical bijection of process collections
\[
\operatorname{Proc}_{\partial}(\mathbb A\otimes\mathbb B,\mathbb C)
\cong
\operatorname{Proc}_{\partial}(\mathbb A,\mathbb B^\bot\parr\mathbb C),
\]
where both sides denote the same underlying Mealy morphisms after the
canonical product reassociations and polarity identifications.
\end{proposition}

\begin{proof}
A process
\[
\mathbb A\otimes\mathbb B\to\mathbb C
\]
is a Mealy morphism
\[
(\mathbb A^+\times\mathbb B^+)\times\mathbb C^-
\rightsquigarrow
\mathbb C^+\times(\mathbb A^-\times\mathbb B^-).
\]
A process
\[
\mathbb A\to\mathbb B^\bot\parr\mathbb C
\]
is a Mealy morphism
\[
\mathbb A^+\times(\mathbb B^\bot\parr\mathbb C)^-
\rightsquigarrow
(\mathbb B^\bot\parr\mathbb C)^+\times\mathbb A^-.
\]
Using
\[
\mathbb B^\bot\parr\mathbb C
=
((\mathbb B^\bot)^\bot\otimes\mathbb C^\bot)^\bot
=
(\mathbb B\otimes\mathbb C^\bot)^\bot,
\]
one identifies
\[
(\mathbb B^\bot\parr\mathbb C)^-
\cong
\mathbb B^+\times\mathbb C^-,
\]
and
\[
(\mathbb B^\bot\parr\mathbb C)^+
\cong
\mathbb B^-\times\mathbb C^+.
\]
Thus the second process type is canonically
\[
\mathbb A^+\times\mathbb B^+\times\mathbb C^-
\rightsquigarrow
\mathbb B^-\times\mathbb C^+\times\mathbb A^-,
\]
which is the same Mealy type as the first, up to product symmetry.
\end{proof}

\section{Multiplicative partial linear syntax}
\label{sec:multiplicative-partial-linear-syntax}

We now define the multiplicative linear syntax interpreted by bidirectional
Mealy processes.  The proof theory is multiplicative linear logic with partial
cut.

\begin{definition}[Formulas]
\label{def:mll-formulas}
Formulas are generated by
\[
A,B
::=
X
\mid
A^\bot
\mid
A\otimes B
\mid
\mathbf 1
\mid
A\parr B
\mid
\bot,
\]
where \(X\) ranges over atomic formulas.
\end{definition}

Linear negation is involutive and satisfies the De Morgan equations
\[
(A^\bot)^\bot=A,
\]
\[
(A\otimes B)^\bot=A^\bot\parr B^\bot,
\]
\[
(A\parr B)^\bot=A^\bot\otimes B^\bot,
\]
\[
\mathbf 1^\bot=\bot,
\qquad
\bot^\bot=\mathbf 1.
\]
Linear implication is treated as an abbreviation:
\[
A\multimap B:=A^\bot\parr B.
\]

\begin{definition}[Contexts]
\label{def:mll-contexts}
A context is a finite list
\[
\Gamma=A_1,\dots,A_n.
\]
The empty context is denoted by
\[
\cdot.
\]
Contexts are considered up to exchange, but not up to weakening or
contraction.
\end{definition}

We use single-conclusion judgments
\[
\Gamma\vdash A.
\]
A derivation of such a judgment will denote a bidirectional Mealy process
\[
\llbracket \Gamma\rrbracket\to \llbracket A\rrbracket.
\]

\section{Interface interpretation of formulas}
\label{sec:interface-semantics-formulas}

Fix an assignment sending each atomic formula \(X\) to a bidirectional Mealy
interface
\[
\llbracket X\rrbracket
=
(\mathbb X^+,\mathbb X^-).
\]

\begin{definition}[Interface interpretation]
\label{def:interface-interpretation}
The interpretation of formulas as bidirectional Mealy interfaces is defined
recursively by:
\[
\llbracket A^\bot\rrbracket
=
\llbracket A\rrbracket^\bot,
\]
\[
\llbracket A\otimes B\rrbracket
=
\llbracket A\rrbracket\otimes\llbracket B\rrbracket,
\]
\[
\llbracket \mathbf 1\rrbracket
=
\mathbf I,
\]
\[
\llbracket A\parr B\rrbracket
=
\left(
\llbracket A^\bot\rrbracket
\otimes
\llbracket B^\bot\rrbracket
\right)^\bot,
\]
and
\[
\llbracket \bot\rrbracket
=
\mathbf I^\bot.
\]
\end{definition}

Since
\[
\mathbf I^\bot=\mathbf I,
\]
the interface interpretations of \(\mathbf 1\) and \(\bot\) have canonically
isomorphic underlying interfaces.  Likewise, at the raw interface level,
\(\parr\) is canonically related to \(\otimes\) by duality.

\begin{definition}[Context interpretation]
\label{def:context-interpretation}
For a context
\[
\Gamma=A_1,\dots,A_n,
\]
define
\[
\llbracket \Gamma\rrbracket
:=
\llbracket A_1\rrbracket\otimes\cdots\otimes
\llbracket A_n\rrbracket.
\]
For the empty context,
\[
\llbracket \cdot\rrbracket=\mathbf I.
\]
\end{definition}

\begin{proposition}[Interpretation respects De Morgan duality]
\label{prop:interface-semantics-demorgan}
The interface interpretation satisfies the linear negation equations up to the
canonical bidirectional interface isomorphisms:
\[
\llbracket (A^\bot)^\bot\rrbracket
\cong
\llbracket A\rrbracket,
\]
\[
\llbracket (A\otimes B)^\bot\rrbracket
\cong
\llbracket A^\bot\parr B^\bot\rrbracket,
\]
\[
\llbracket (A\parr B)^\bot\rrbracket
\cong
\llbracket A^\bot\otimes B^\bot\rrbracket,
\]
and
\[
\llbracket \mathbf 1^\bot\rrbracket
\cong
\llbracket \bot\rrbracket.
\]
\end{proposition}

\begin{proof}
This follows directly from the definitions of formula interpretation, duality,
tensor, and par on bidirectional interfaces.
\end{proof}

\section{Proof rules and process interpretation}
\label{sec:proof-rules-process-interpretation}

A derivation
\[
\pi:\Gamma\vdash A
\]
is interpreted as a bidirectional Mealy process
\[
\llbracket \pi\rrbracket:
\llbracket \Gamma\rrbracket
\to
\llbracket A\rrbracket.
\]

We list the basic rules and their interpretations.

\subsection{Identity}

The identity rule is
\[
\mathrm{infer}[\mathrm{id}]
      {A\vdash A}
      {}
\]
and is interpreted by the identity process
\[
\llbracket \mathrm{id}_A\rrbracket
=
1_{\llbracket A\rrbracket}.
\]

\subsection{Exchange}

Exchange permutes the order of formulas in a context.  It is interpreted by
the symmetry isomorphisms of the tensor product of bidirectional interfaces.

\subsection{Tensor introduction}

The tensor introduction rule is
\[
\mathrm{infer}[\otimes I]
      {\Gamma,\Delta\vdash A\otimes B}
      {\Gamma\vdash A \& \Delta\vdash B}.
\]
If
\[
\llbracket \pi\rrbracket:
\llbracket\Gamma\rrbracket\to\llbracket A\rrbracket
\]
and
\[
\llbracket \rho\rrbracket:
\llbracket\Delta\rrbracket\to\llbracket B\rrbracket,
\]
then
\[
\llbracket \pi\otimes\rho\rrbracket
=
\llbracket \pi\rrbracket\otimes\llbracket \rho\rrbracket,
\]
after the canonical context reassociations.

\subsection{Tensor unit}

The tensor unit introduction rule is
\[
\mathrm{infer}[\mathbf 1 I]
      {\cdot\vdash\mathbf 1}
      {}
\]
and is interpreted by the identity process on the unit interface
\[
1_{\mathbf I}:\mathbf I\to\mathbf I.
\]

\subsection{Linear implication}

Linear implication is the abbreviation
\[
A\multimap B:=A^\bot\parr B.
\]
The introduction rule is
\[
\mathrm{infer}[\multimap I]
      {\Gamma\vdash A\multimap B}
      {\Gamma,A\vdash B}.
\]
It is interpreted by the canonical bidirectional currying isomorphism
\[
\operatorname{Proc}_{\partial}
(\llbracket\Gamma\rrbracket\otimes\llbracket A\rrbracket,
 \llbracket B\rrbracket)
\cong
\operatorname{Proc}_{\partial}
(\llbracket\Gamma\rrbracket,
 \llbracket A\multimap B\rrbracket),
\]
from Proposition~\ref{prop:bidirectional-currying}.  This isomorphism is
induced by polarity reversal and product reassociation; it is not an additional
closed-structure operation.

The elimination rule is
\[
\mathrm{infer}[\multimap E]
      {\Gamma,\Delta\vdash B}
      {\Gamma\vdash A\multimap B \& \Delta\vdash A}.
\]
It is interpreted by partial GoI composition, after the canonical
reassociations identifying the function process with a process consuming an
\(A\)-input.  The rule is defined precisely when the resulting feedback
problem is trace-class.

\subsection{Cut}

The cut rule is
\[
\mathrm{infer}[\mathrm{cut}]
      {\Gamma,\Delta\vdash B}
      {\Gamma\vdash A \& \Delta,A\vdash B}.
\]
It is a partial rule.  If
\[
\llbracket\pi\rrbracket:
\llbracket\Gamma\rrbracket\to\llbracket A\rrbracket
\]
and
\[
\llbracket\rho\rrbracket:
\llbracket\Delta\rrbracket\otimes\llbracket A\rrbracket
\to
\llbracket B\rrbracket,
\]
then, after the canonical associativity and symmetry identifications,
\[
\llbracket \operatorname{cut}(\pi,\rho)\rrbracket
=
\llbracket\rho\rrbracket
\circ_{\partial}
(\llbracket\pi\rrbracket\otimes 1_{\llbracket\Delta\rrbracket}),
\]
provided this partial composite is defined.

Thus cut is admissible precisely when the corresponding open pipeline is
trace-class.

\section{Soundness of multiplicative partial linear logic}
\label{sec:soundness-multiplicative-partial-linear-logic}

\begin{definition}[Defined derivation]
\label{def:defined-derivation}
A derivation is \textbf{defined} if every partial composition required in its
process interpretation is trace-class.
\end{definition}

\begin{theorem}[Soundness]
\label{thm:soundness-partial-mll}
Every defined derivation
\[
\pi:\Gamma\vdash A
\]
denotes a bidirectional Mealy process
\[
\llbracket\pi\rrbracket:
\llbracket\Gamma\rrbracket\to\llbracket A\rrbracket
\]
in
\[
\operatorname{Int}_{\partial}(\mathsf{Mly}(\mathcal E)).
\]
The interpretation of cut is partial GoI composition, and is defined exactly
when the associated open pipeline is trace-class.
\end{theorem}

\begin{proof}
The proof is by induction on derivations.

The identity rule is interpreted by the identity process.  Exchange is
interpreted by the symmetry of the tensor product.  Tensor introduction is
interpreted by tensoring processes, which is well typed by
Proposition~\ref{prop:tensor-processes-well-typed}.  The tensor unit rule is
interpreted by the identity process on the unit interface.

The implication introduction rule is interpreted by the bidirectional currying
isomorphism of Proposition~\ref{prop:bidirectional-currying}.  Implication
elimination and cut are interpreted by partial GoI composition.  By definition
of a defined derivation, the necessary partial composites are trace-class.
Hence the interpretation exists as a bidirectional Mealy process.

Therefore every defined derivation has a well-defined process interpretation
of the required type.
\end{proof}

\section{Cut execution}
\label{sec:cut-execution}

The semantic meaning of cut is not merely that two processes are composed.
Cut is executed by solving the hidden feedback equation of the associated GoI
pipeline.

\begin{theorem}[Cut execution]
\label{thm:cut-execution}
Let
\[
\pi:\Gamma\vdash A
\]
and
\[
\rho:\Delta,A\vdash B
\]
be derivations.  Suppose the denotational cut
\[
\llbracket\rho\rrbracket
\circ_{\partial}
(\llbracket\pi\rrbracket\otimes 1_{\llbracket\Delta\rrbracket})
\]
is defined.  Then the denotation of the cut is computed by solving the unique
feedback equation of the corresponding open pipeline and substituting the
solution into the transition and output maps of the two processes.
\end{theorem}

\begin{proof}
Let
\[
F=
\llbracket\pi\rrbracket\otimes 1_{\llbracket\Delta\rrbracket}
\]
and
\[
G=
\llbracket\rho\rrbracket.
\]
The denotational cut is
\[
G\circ_{\partial}F.
\]
By assumption, this partial composite is defined.  Hence, by
Theorem~\ref{thm:goi-execution-criterion}, the feedback witness projection
\[
\pi_{G,F}:W_{G,F}\to D_{G,F}
\]
is an isomorphism.  Therefore the feedback equation
\[
h
=
\omega_G^-
\left(
\left(
\omega_F^+((a^+,h),x),
c^-
\right),
y
\right)
\]
has a unique internal solution
\[
h_{G,F}(a^+,c^-,x,y).
\]
The component execution formula of
Theorem~\ref{thm:component-execution-formula} says that the transition and
output of the cut are obtained by substituting this solution into the open
pipeline.  This gives the claimed computation of the cut.
\end{proof}

\begin{remark}[Semantic cut execution]
\label{rem:semantic-cut-execution}
The theorem is a semantic cut-execution result.  It does not assert a
syntactic normalization theorem for a proof-rewriting system.  It says that
whenever a cut is denotationally defined, its interpretation is computed by
Geometry-of-Interaction feedback.
\end{remark}

\section{Structural laws for partial cut}
\label{sec:structural-laws-partial-cut}

The paracategory laws give the structural equations of partial cut.

\begin{proposition}[Identity cut]
\label{prop:identity-cut}
For every process
\[
F:\mathbb A\to\mathbb B,
\]
one has
\[
F\circ_{\partial}1_{\mathbb A}=F,
\qquad
1_{\mathbb B}\circ_{\partial}F=F,
\]
on the trace-class domains of the displayed composites.
\end{proposition}

\begin{proof}
These are the yanking equations for the partial trace inherited from the
previous chapter.  Composing with a wire process forms an open pipeline whose
feedback equation is the wire equation, and the traced result is canonically
identified with \(F\).
\end{proof}

\begin{proposition}[Associativity of partial cut]
\label{prop:associativity-partial-cut}
Let
\[
F:\mathbb A\to\mathbb B,
\qquad
G:\mathbb B\to\mathbb C,
\qquad
H:\mathbb C\to\mathbb D
\]
be bidirectional Mealy processes.  Whenever the two bracketed composites are
defined, one has
\[
H\circ_{\partial}(G\circ_{\partial}F)
=
(H\circ_{\partial}G)\circ_{\partial}F.
\]
\end{proposition}

\begin{proof}
Both sides are obtained by forming the same three-stage open pipeline and
closing the same hidden counterflow interfaces.  The two bracketings correspond
to tracing these interfaces in different orders.  The staged-feedback theorem
from the previous chapter identifies the iterated traces with the simultaneous
trace on their common trace-class domain.  Hence the two composites agree.
\end{proof}

\begin{proposition}[Tensor compatibility]
\label{prop:tensor-compatibility-partial-cut}
Let
\[
F_i:\mathbb A_i\to\mathbb B_i,
\qquad
G_i:\mathbb B_i\to\mathbb C_i
\]
for \(i=1,2\).  Whenever the relevant partial composites are defined, one has
\[
(G_1\circ_{\partial}F_1)\otimes(G_2\circ_{\partial}F_2)
\cong
(G_1\otimes G_2)\circ_{\partial}(F_1\otimes F_2).
\]
\end{proposition}

\begin{proof}
The open pipeline for
\[
F_1\otimes F_2
\]
followed by
\[
G_1\otimes G_2
\]
is canonically the tensor of the two open pipelines for \((F_1,G_1)\) and
\((F_2,G_2)\).  The superposing law for the partial trace identifies the trace
of the tensor with the tensor of the traces on the common trace-class domain.
\end{proof}

\section{Partial linear logic of Mealy GoI}
\label{sec:partial-linear-logic-mealy-goi}

\begin{theorem}[Partial linear logic of Mealy GoI]
\label{thm:partial-linear-logic-mealy-goi}
Let \(\mathcal E\) have finite limits.  The partial Int construction applied
to the partially traced symmetric monoidal category
\[
\mathsf{Mly}(\mathcal E)
\]
yields a symmetric monoidal paracategory
\[
\operatorname{Int}_{\partial}(\mathsf{Mly}(\mathcal E))
\]
with involutive duality.

This structure interprets multiplicative linear logic with partial cut:
\begin{enumerate}
\item formulas are interpreted as bidirectional Mealy interfaces;
\item proofs are interpreted as bidirectional Mealy processes;
\item tensor is interpreted by product of interfaces and parallel composition
of processes;
\item linear negation is interpreted by reversal of polarity;
\item par is interpreted by De Morgan duality;
\item cut is interpreted by partial GoI composition.
\end{enumerate}
A cut is defined precisely when the associated open pipeline is trace-class,
equivalently when its hidden feedback equation has a unique internal solution.
When defined, the denotation of the cut is computed by substituting that unique
solution into the open-pipeline transition and output maps.
\end{theorem}

\begin{proof}
The existence of the partial trace on
\[
\mathsf{Mly}(\mathcal E)
\]
is the main result of the previous chapter.  The construction of
\[
\operatorname{Int}_{\partial}(\mathsf{Mly}(\mathcal E))
\]
and its symmetric monoidal paracategory structure is
Theorem~\ref{thm:partial-int-paracategory}.  The involutive duality is given
by reversal of bidirectional polarity.

The interpretation of formulas and contexts is
Definition~\ref{def:interface-interpretation} and
Definition~\ref{def:context-interpretation}.  The interpretation of proof
rules is given in Section~\ref{sec:proof-rules-process-interpretation}.
Soundness for defined derivations is Theorem~\ref{thm:soundness-partial-mll}.

The characterization of defined cut as trace-class feedback is
Theorem~\ref{thm:goi-execution-criterion}.  The component computation of a
defined cut by substituting the unique feedback solution is
Theorem~\ref{thm:component-execution-formula} and
Theorem~\ref{thm:cut-execution}.
\end{proof}

\section{Conceptual summary}
\label{sec:goi-conceptual-summary}

The construction of this chapter may be summarized as follows.

First, a bidirectional Mealy process
\[
F:\mathbb A\to\mathbb B
\]
is a Mealy morphism
\[
F:
\mathbb A^+\times\mathbb B^-
\rightsquigarrow
\mathbb B^+\times\mathbb A^-.
\]
It carries positive behaviour forward and negative behaviour backward.

Second, the composite of
\[
F:\mathbb A\to\mathbb B
\]
and
\[
G:\mathbb B\to\mathbb C
\]
is not ordinary sequential composition.  It is Geometry-of-Interaction
execution:
\[
G\circ_{\partial}F
=
\operatorname{Tr}^{\mathbb B^-}(H_{G,F}).
\]
Thus composition is defined precisely when the hidden counterflow equation has
a unique internal solution.

Third, the proof-theoretic cut rule is interpreted by this same operation.  A
cut is admissible exactly when the corresponding feedback equation is
trace-class.

Finally, when a cut is defined, its denotation is computed by solving the
feedback equation
\[
h
=
\omega_G^-
\left(
\left(
\omega_F^+((a^+,h),x),
c^-
\right),
y
\right)
\]
and substituting the unique solution into the open-pipeline transition and
output maps.

Thus the chapter identifies the chain
\[
\boxed{
\text{partial trace}
\quad\Longrightarrow\quad
\text{partial GoI execution}
\quad\Longrightarrow\quad
\text{partial cut}
\quad\Longrightarrow\quad
\text{multiplicative linear process logic}.
}
\]